\title{Large Language Models Can Automatically Engineer Features for Few-Shot Tabular Learning}

\begin{abstract}
Large Language Models (LLMs) possess an extraordinary capacity to address complex and novel reasoning challenges, positioning them as promising tools for tabular learning—a foundational component across numerous real-world scenarios. In this study, we introduce \textsf{FeatLLM}, an innovative in-context learning approach wherein LLMs serve as automatic feature constructors, producing feature-enriched datasets optimized for tabular prediction tasks. These synthesized features subsequently feed into a straightforward downstream predictor, such as linear regression, facilitating strong few-shot performance. Critically, our \textsf{FeatLLM} framework confines the inference process to the utilization of the simple downstream predictive model alongside the discovered features, obviating the requirement to interact with the LLM during inference. This is in contrast with other LLM-based techniques that necessitate repeated querying of the model at test time. Notably, \textsf{FeatLLM} operates through a mere API-level interface with LLMs, sidestepping both prompt length restrictions and inference-time costs. Empirical results spanning diverse tabular benchmarks substantiate that \textsf{FeatLLM} reliably generates superior predictive rules, yielding, on average, a 10\% improvement over contemporaries such as TabLLM and STUNT.
\end{abstract}

\section{Introduction}
The emergence of Large Language Models (LLMs) in recent years has unveiled their substantial proficiency at solving unfamiliar tasks through natural language responses, typically without recourse to bespoke fine-tuning for specific tasks~\cite{creswell2022selection,imani2023mathprompter,taylor2022galactica}. LLMs are trained on vast textual resources, granting them a wealth of intrinsic knowledge that can act as a proxy for human expert insight across various disciplines~\cite{Genomes2021,singhal2023large,wilcox2003role}, and are pivotal to automation efforts in numerous applications such as conversational analysis~\cite{finch2023leveraging} and program repair~\cite{fan2023automated}. In particular, by embedding concise task descriptions and a handful of labeled samples within their prompts, LLMs have demonstrated the ability to intuitively capture task context and identify relevant features and exemplars~\cite{brown2020language,zhao2023survey}. Such findings underscore the capability of LLMs to act as proficient feature engineers through data interpretation.

The practice of leveraging the knowledge and inference power of LLMs has extended to tabular learning contexts. For example, domain-relevant knowledge, such as the association between advanced age and higher disease prevalence, can enhance predictive modeling in healthcare. Recent studies have operationalized this by formalizing tasks and feature explanations in natural language, converting structured data into LLM-ingestible formats, and submitting these to LLMs for processing~\cite{dinh2022lift,wang2023anypredict}. Follow-up strategies often involve in-context learning~\cite{nam2023semi} or efficient, limited-parameter adaptation techniques~\cite{dinh2022lift,hegselmann2023tabllm}. In few-shot settings, the knowledge embedded in LLMs has enabled performance gains over conventional tabular learning approaches~\cite{hegselmann2023tabllm}.

Existing approaches leveraging LLMs for tabular data analysis face several notable challenges. End-to-end predictive models necessitate at least one LLM evaluation per input instance, which results in substantial computational overhead. Achieving high predictive accuracy often depends on fine-tuning the LLM, a process that is generally unworkable when dealing with LLMs lacking public training pipelines or accessible fine-tuning interfaces. This issue is particularly pressing given that many recent state-of-the-art LLMs only allow users to interact with them through limited inference APIs~\cite{anil2023palm,openai2023gpt4}. Moreover, these methods typically perform poorly with extended input prompts~\cite{liu2023lost}; as the tabular dataset’s feature count rises and text serialization expands, method efficacy decreases or becomes impractical. Collectively, these limitations hinder the adoption of LLM-based tabular solutions in real-world settings.

To address these obstacles, our method departs from the conventional framework of LLMs generating end-to-end predictions, positioning them instead as instruments for feature engineering. This alternative harnesses their extensive pretrained knowledge more effectively. We introduce an innovative in-context learning paradigm, \textsf{FeatLLM}, designed to elucidate the “criteria” that guide LLM-generated predictions. For example, in diagnosing diseases, the LLM can deduce and articulate explicit rules identifying which combinations of features correspond to the diagnosis. These generated criteria (or rules) serve as the foundation for new, synthesized features that can replace or supplement existing ones for downstream predictive tasks. This approach transitions the LLM’s primary function from direct inference to feature construction, enabling the use of straightforward predictive models downstream and markedly reducing inference time after feature extraction. Utilizing the in-context learning capabilities of LLMs allows us to develop features by incorporating illustrative examples within prompts, eliminating the need for fine-tuning. Additionally, we address issues related to large prompt sizes by incorporating ensemble learning techniques, such as feature bagging~\cite{breiman2001random}, to enhance scalability.

\textsf{FeatLLM} decomposes the prompt for feature engineering into two key sub-tasks: (1) comprehending the problem context while inferring potential dependencies between features and targets; and (2) leveraging insights from the initial analysis, together with few-shot examples, to formulate discriminative rules that differentiate between classes. This deliberate reasoning framework guides LLMs to disregard non-informative features, thus improving the quality of extracted rules. Once generated, these rules are executed on the data (by parsing them as code), and each sample is mapped to a binary feature reflecting compliance with each rule. Subsequently, a simple model is trained on these constructed features to estimate the likelihood of each class. To enhance robustness, an ensemble scheme is used, where feature generation and feature bagging are repeated multiple times. By strategically varying the count of features and the selection of samples in each bagging iteration, the difficulties posed by excessively long prompts are alleviated. \textsf{FeatLLM} operates independently of training demands and incurs minimal inference overhead, making LLM-based feature engineering broadly accessible for a range of practical tasks.

Our empirical assessment of \textsf{FeatLLM} spans 13 diverse tabular datasets in low-data settings, demonstrating its consistently strong and stable results. We find that LLMs effectively integrate both external knowledge and information extracted from the provided data to create robust features. Our proposed method consistently surpasses modern few-shot learning baselines in various experimental conditions. The corresponding source code is publicly accessible via the anonymized GitHub repository at~\texttt{https://github.com/Sungwon-Han/FeatLLM}.

\section{Related Work}

\subsection{Few-Shot Learning with Tabular Data} Progress in few-shot learning has made it feasible to derive broadly applicable data representations while incurring low annotation costs~\cite{chen2018closer}. The field initially concentrated on image-based problems, but few-shot paradigms have since shown their versatility in handling diverse modalities, encompassing text, audio, and, more recently, tabular data~\cite{majumder2022few,nam2023stunt,schick2021exploiting}. Nevertheless, learning effective models from only a handful of labeled examples remains susceptible to the pitfalls of spurious patterns~\cite{bartlett2021deep}. To address this challenge, contemporary research has drawn on auxiliary information and infused inductive biases relevant to specific domains by leveraging prior knowledge. Some strategies augment the training regime with large quantities of unlabeled data through well-chosen augmentation strategies for semi-supervised learning~\cite{ucar2021subtab,yoon2020vime}; others seek robust parameter initialization via unsupervised, task-agnostic pre-training~\cite{somepalli2021saint}. Typical unsupervised techniques involve objectives such as masking or introducing corruption to parts of the input, followed by reconstruction tasks~\cite{arik2021tabnet,majmundar2022met}, or leveraging data augmentation to support contrastive learning schemes~\cite{bahri2022scarf,chen2020simple}. Despite these advances, the inherent heterogeneity of tabular data presents challenges in devising general augmentation techniques, which has, in turn, limited the success of such methods in few-shot contexts~\cite{nam2023stunt}.

A different avenue has explored training models across a heterogeneous collection of tabular datasets—including those that differ from the intended application—and then using transfer learning techniques to adapt to specific target tasks~\cite{zhang2023generative,levin2022transfer,zhu2023xtab}. One such method, TransTab~\cite{wang2022transtab}, leverages transformer architectures to capture inter-column semantic relationships primarily via column names, beyond solely relying on table content. The extraction of knowledge at the level of columns and tables enables the system to perform zero-shot or few-shot predictions on entirely novel tabular problems. Similarly, TabPFN~\cite{hollmann2022tabpfn} combines various distributional types to synthetically generate data for pre-training a Bayesian neural network; after this phase, the model can make predictions for new tasks by consuming both the training and testing samples directly, bypassing additional training. Leveraging prior knowledge distilled from numerous datasets has yielded performance that not only surpasses conventional tree-based methods but also demonstrates tangible benefits in few-shot scenarios.

\subsection{Language-Interfaced Tabular Learning} An additional promising direction involves the integration of large language models (LLMs)~\cite{anil2023palm,openai2023gpt4}, which are trained on vast and varied textual resources and exhibit strong generalization even to previously unseen tasks~\cite{brown2020language}. Contemporary research has focused on utilizing the extensive stored knowledge within LLMs for learning from tabular data. Common practices entail linearizing tabular records into text and incorporating them as part of the prompts provided to LLMs~\cite{dinh2022lift,hegselmann2023tabllm,wang2023anypredict}. Prompts often include raw tabular entries, feature descriptions, and details on the tasks to be performed, enabling LLMs to reason over both the dataset structure and task requirements. Developments in this domain range from targeted model fine-tuning via parameter-efficient methods such as LoRA~\cite{hu2021lora} or IA3~\cite{liu2022few}, to in-context learning approaches where a set of exemplar demonstrations is simply inserted into the prompts, obviating the need for further tuning or retraining~\cite{dinh2022lift,hegselmann2023tabllm,nam2023semi,slack2023tablet}.

Although previous methods have demonstrated efficacy in enhancing few-shot tabular learning, their reliance on passing every inference through the LLM introduces hurdles for deployment within industrial environments. This work seeks to shift away from the traditional black-box paradigm, in which predictions for individual samples are produced end-to-end by an LLM. Rather than relying on the LLM for direct prediction, the approach proposed here leverages the LLM to extract class-specific rules or conditions. \textsf{FeatLLM} transforms these inferred rules into executable code that generates novel features, thereby supporting inference independent of the LLM, while utilizing in-context learning for rule induction based on both background knowledge and limited labeled examples.

\section{Methods}
The central objective of \textsf{FeatLLM} is to discover the set of “rules,” or defining conditions for each output class, using the limited annotated samples as demonstrated in Figure~\ref{fig:main_model}. These rules are first synthesized by the LLM, after which they are converted into programmatic representations that serve as binary features for each instance; a value indicates if a particular rule applies to the instance. These binary feature vectors are then combined through a dot product operation with non-negative parameters, yielding class probability estimates. During training, this model learns the significance of each rule-feature in predicting class membership (see Section~\ref{Sec:training}). The entire process utilizes repeated bagging for robust feature aggregation, followed by ensembling for improved prediction. The subsequent sections detail the overall problem setting and elucidate the procedures at each stage.

\vspace*{-0.12in}
\paragraph{Problem formulation.} Consider a tabular dataset composed of $N$ labeled examples, represented as $\mathcal{D} = \{(\mathbf{x}^i, \mathbf{y}^i)\}_{i=1}^{N}$. Each data point $\mathbf{x}^i$ is a $d$-dimensional vector corresponding to $d$ distinct input features, and each feature has an associated name in natural language, denoted $F = \{f_j\}_{j=1}^{d}$. Definitions for these feature names are specified separately. We focus on a classification framework where each label $\mathbf{y}^i$ takes a value from a specified set of categories. During $k$-shot learning experiments, a subset of $k$ samples (with $k < N$) is selected at random from the labeled data to train the model.

\subsection{Prompt Design for Extracting Rules}
\label{Sec:prompt_design}
To facilitate rule extraction by LLMs through effective reasoning, we structure the problem-solving sequence in a way that reflects the systematic approach typically adopted by domain experts in tabular analysis. The prompt is organized around three primary components (see Fig.~\ref{fig:prompt_for_rule} and Appendix~\ref{sec:example_rule}):

\vspace*{-0.12in}
\paragraph{Basic information description.} This component supplies key contextual details required to address the problem at hand (indicated by orange text in Fig.~\ref{fig:prompt_for_rule}). It encompasses a task description specifying label information, feature descriptions, and example instances. The nature of the task is stated as a question (for instance: ``Does this patient's myocardial infarction complications data indicate chronic heart failure? Yes or no?"; refer to Appendix~\ref{sec:task_desc} for more examples). Each feature description characterizes its value type and may optionally provide the feature's definition; for categorical variables, illustrative category values can be listed. To serve as exemplars, select training samples are transformed into textual form together with their corresponding true labels, using the serialization approach adopted in previous literature~\cite{dinh2022lift,hegselmann2023tabllm}:

\begin{align}
&\text{Serialize}(\mathbf{x}^i,\mathbf{y}^i, F) = \nonumber \\ 
&``\ f_1\ \text{is}\ \mathbf{x}^i_1.\ ... \ f_d\  \text{is}\  \mathbf{x}^i_d.\ \ \text{Answer:}\  \mathbf{y}^i\ ”
\end{align}

\vspace*{-0.12in}
\paragraph{Reasoning instruction.} Our goal is to strengthen the LLM’s reasoning abilities by introducing explicit reasoning directives (refer to the blue-highlighted section in Fig.~\ref{fig:prompt_for_rule}). The provided reasoning instruction opens with a statement akin to those used in the chain-of-thought paradigm~\cite{wei2022chain}. Subsequently, the LLM’s inference workflow is decomposed into two distinct stages. During the initial stage, the model is prompted to ascertain the causal dynamics or general associations between the features and the provided task description, drawing only upon its pre-existing knowledge base or common sense rather than on example demonstrations. This facilitates a thorough retrieval of the LLM’s background knowledge about the problem context. In the following stage, the model integrates information from both the first stage and sample demonstrations to distill classification rules pertinent to each class. Implementing this bifurcated reasoning sequence helps minimize the chance of detecting irrelevant correlations in uninformative columns and directs the model's attention toward essential features. To avoid issues such as extracting an insufficient or excessive number of rules—which could lead to underfitting or an overly lengthy prompt—we impose a fixed count of rules to extract for every class (specifically, 10)\footnote{Further discussion on the choice of rule count is found in Section~\ref{sec:analysis}.}. Furthermore, when deriving rules, the LLM references the type information for each feature, as described in the basic information segment, thereby reducing potential type mismatches within the formulated rules.

\vspace*{-0.12in}
\paragraph{Response instruction.} To promote consistent and easily interpretable outputs for the deduced rules, we provide the LLM with structured response directives (see the yellow-highlighted text in Fig.~\ref{fig:prompt_for_rule}). The prompt details explicit requirements concerning the formatting of responses assigned to each answer class (that is, the structure for answers), as well as the template that each rule is to follow (the rule structure). Through these directives, the LLM is able to select the appropriate format according to the feature type in question. Additionally, barring extra constraints, the LLM is permitted to logically combine individual rules—using connectors such as AND or OR. An example of the resulting output can be found in Appendix~\ref{sec:outcome-rule}.

\subsection{Inferring Class Likelihood via Rules}
\label{Sec:training}

\paragraph{Parsing rules for feature generation.} In this phase, the previously generated rules are converted into additional binary features. These features are produced for every class and specify whether a sample complies with the corresponding rules for that class (i.e., values of '0' or '1'). They serve as indicators to help estimate the probability that a sample falls into a particular class. However, since these rules are articulated in natural language by the LLM, an effective parsing mechanism is necessary to enable their automatic transformation into features. Although we impose structured response guidelines to facilitate easier parsing, outputs from the LLM can still exhibit syntactic inconsistencies~\cite{kovavc2023large}. For example, rather than using symbols like ``$>$” or ``$<$”, the LLM might opt for equivalent phrases such as ``greater than" or ``less than." Such inconsistencies complicate the parsing process.

Rather than constructing elaborate code to parse potentially noisy language, we instead employ the LLM to aid in parsing. The LLM is adept at grasping semantic content even in the presence of varying syntax, and it is proficient in generating concise code according to prompts, owing to its code-based training data. To capitalize on these strengths, our prompt includes the target function name, details of inputs and outputs, and the extracted rules, before submitting it to the LLM. Illustrative prompts and their resulting outputs can be found in Figures~\ref{fig:prompt_for_parsing} and~\ref{sec:example_parsing}-\ref{sec:example_parse_outcome} in the Appendix. The code output by the LLM is subsequently run using Python’s \texttt{exec()} function to process and convert the data.

\vspace*{-0.12in}
\paragraph{Inferring class likelihood.} Once binary rule-based features have been established for each class, a straightforward approach for estimating the likelihood of a sample belonging to a specific class involves tallying the number of rules that are satisfied (i.e., computing the sum of binary features for each class). Nevertheless, some rules or features contribute more strongly to class discrimination than others, so it is important to determine their relative impact via learning from labeled data. To this end, we deploy a standard ML model to assign weights to each class's binary features. For instance, a linear model without an intercept can be utilized. We represent the collection of binary features generated from sample $\mathbf{x}^i$ for class $k$ as $\mathbf{z}_k^i \in \{0, 1\}^R$, where $R$ refers to the count of rules per class, and $c$ denotes the total number of classes. The class probability vector $\mathbf{p}^i$ is then determined as follows:

\begin{align}
\text{logit}_k^i &= \max(\mathbf{w}_k, 0) \cdot \mathbf{z}_k^i, \nonumber \\
\mathbf{p}^i &= \text{Softmax}({[\text{logit}_{1}^i, ..., \text{logit}_{c}^i]}). \label{eq:infer_prob}
\end{align}

In this context, $\mathbf{w}_k$ denotes the tunable parameters of the linear classifier, reflecting the relative contribution of each feature. By constraining these weights to reside within a positive subspace, we guarantee that features proposed by the LLM cannot detract from a class’s probability during learning. Afterward, the logit outputs are mapped to class probabilities via the softmax transformation. The cross-entropy loss is minimized to fit the weights $\mathbf{w}_k$, leveraging only a small set of training instances.

\vspace*{-0.12in}
\paragraph{Ensembling with bagging.} To conclude, we repeatedly iterate the entire learning process, producing an ensemble of inference models whose collective prediction is based on the aggregation of individual models. Specifically, for $T$ distinct runs producing trained weights $\{\mathbf{w}[t]\}_{t=1}^T$, final class probabilities are calculated as the average over $\mathbf{p}^i$ derived for each $\mathbf{w}[t]$, consistent with Eq.~\ref{eq:infer_prob}.

To inject diversity into the ensemble, we raise the temperature parameter during LLM inference, amplifying stochasticity in rule creation. Furthermore, since the sequence of few-shot samples can influence the LLM’s generation~\cite{lu2022fantastically}, we randomly permute the demonstration order; an analysis of the resulting variety appears in Appendix~\ref{sec:diversity}. To further promote heterogeneity in the ensemble, particularly as training set sizes or feature counts grow, we employ bagging, selecting random subsets of either features or samples in each trial. This ensemble strategy confers two central benefits. First, it enables the system to recover from inaccuracies by the LLM in any particular run, as successful rule generations across other runs enhance framework stability; this relates to LLM self-consistency~\cite{wang2022self}, whereby rules consistently produced are likelier correct. Second, the ensemble mitigates input prompt size constraints inherent to LLMs: when tables grow too large for a single prompt, bagging naturally curtails input size per trial. While individual models may not represent all feature associations, assembling multiple weak learners across diverse trials compensates for any missing coverage in a single pass.

\section{Experiments}
We assess the effectiveness of \textsf{FeatLLM} across a diverse suite of tabular datasets, supplementing our results with analyses to deepen understanding of the framework's operational dynamics. Owing to space limitations, extended experimental results—including tests involving alternative LLMs, detailed qualitative studies, and more—are available in the Appendix. 

\subsection{Performance Evaluation}
\paragraph{Datasets.}
Our empirical study is based on 13 datasets involving binary or multi-class classification tasks: (1) Adult~\cite{asuncion2007uci}, designed to estimate whether an individual’s income surpasses \$50,000 per year; (2) Bank~\cite{moro2014data}, targeting the prediction of customer subscription to term deposits; (3) Blood~\cite{yeh2009knowledge}, concerned with the likelihood of donor return for future donations; (4) Car~\cite{kadra2021well}, evaluating automotive quality; (5) Communities~\cite{misc_communities_and_crime_183}, focused on forecasting crime rates within defined regions; (6) Credit-g~\cite{kadra2021well}, classifying creditworthiness as either good or bad; (7) Diabetes\footnote{\texttt{kaggle.com/datasets/uciml/pima-indians-diabetes-database}}, predicting diabetes onset for a patient; (8) Heart\footnote{\texttt{kaggle.com/datasets/fedesoriano/heart-failure-prediction}}, concerning the identification of coronary artery disease; (9) Myocardial~\cite{misc_myocardial_infarction_complications_579}, aiming to detect chronic heart failure. 

To broaden the scope, we introduce novel and synthetic benchmarks generally overlooked in LLM pre-training: (10) Cultivars, for forecasting the expected grain yield of soybean cultivars\footnote{\texttt{archive.ics.uci.edu/dataset/913}}; (11) NHANES, to predict an individual's age group based on health and demographic data\footnote{\texttt{archive.ics.uci.edu/dataset/887}}; (12) Sequence-type, classifying sequences as arithmetic, geometric, Fibonacci, or Collatz; and (13) Solution-mix, determining whether the concentration of a four-solution mixture is above 0.5. Cultivars and NHANES, having only recently been included in the UCI Machine Learning Repository (post-September 2023), were unavailable during LLM API (gpt-3.5-turbo-0613) pre-training and hence were not contained in its knowledge base. The final two datasets are synthetic, thereby precluding memorization by the LLM and ensuring reliable evaluation.

Details on data size and attribute complexity are provided in Table~\ref{Tab:basic-desc}.
Each dataset includes unambiguous attribute names and accompanying descriptions. A few, like `Communities' and `Myocardial,' feature over 100 attributes, introducing substantial complexity and posing challenges for LLMs with restricted context windows. 

\vspace*{-0.12in}
\paragraph{Baselines.}
Nine baseline methods are incorporated for comparative evaluation. The initial group comprises classic tabular learning models: (1) Logistic regression (LogReg), (2) XGBoost~\cite{chen2016xgboost}, and (3) RandomForest~\cite{ho1995random}. The following two leverage large quantities of unlabeled data: (4) SCARF~\cite{bahri2022scarf} and (5) STUNT~\cite{nam2023stunt}, although, in practice, amassing sufficient unlabeled samples may prove impractical. We also include a modern synthetic-data-based baseline: (6) TabPFN~\cite{hollmann2022tabpfn}, which relies on pretraining over artificially generated datasets with varied distributions. The remaining baselines apply LLM-based methods: (7) In-context learning (In-context)~\cite{wei2022emergent}, (8) TABLET~\cite{slack2023tablet}, and (9) TabLLM~\cite{hegselmann2023tabllm}. In-context learning places few-shot labeled instances directly within the prompt for inference without further optimization. TABLET augments prompt information by incorporating hints produced by an external classifier in the form of rule sets and prototypes to improve decision accuracy. Finally, TabLLM leverages parameter-efficient techniques such as IA3~\cite{liu2022few} for finetuning the LLM using limited training samples.

\vspace*{-0.12in}
\paragraph{Implementation details.}
\textsf{FeatLLM} is implemented with GPT-3.5 serving as the foundational LLM, though the framework is not restricted to any specific LLM and can accommodate various options (see Appendix~\ref{sec:backbones} for experiments with alternative LLM backbones). The LLM inference process employs a temperature of 0.5, while the top-p sampling remains at the API’s default value of 1. We configure the ensemble size to 20 and extract 10 rules. Further analysis of how different hyper-parameters influence performance is available in Figure~\ref{fig:hyperparameter2} and Appendix~\ref{sec:hyper-parameter}.

For the linear model, optimization is handled via the Adam algorithm with a learning rate of 0.01, trained over 200 epochs. Optimal epoch determination relies on $k$-fold cross-validation. When reproducing baselines, we apply GPT-3.5 for in-context learning and T0~\cite{sanh2022multitask} in TabLLM, maintaining consistency with its original protocol. Note that TabLLM restricts model selection to LLMs with open-access checkpoints, so GPT-3.5 cannot be utilized within this framework. Classical machine learning algorithms such as LogReg, XGBoost, and RandomForest are tuned by means of Grid-search in combination with $k$-fold cross-validation. Depending on the dataset, $k$ is set to 2 or 4, always ensuring each class is represented in the training partitions. Further implementation specifics can be found in Appendix~\ref{sec:baselines}.

\vspace*{-0.12in}
\paragraph{Results.}
Tables~\ref{tab:main1} and \ref{tab:main2} display comparative results across multiple datasets. Each experiment is run three times using different random splits for the training and test sets; the reported AUC values reflect the mean and standard deviation over these runs. \textsf{FeatLLM} regularly appears as either the best or second-best model in relation to various baseline methods. Remarkably, our approach attains a similar level of accuracy with only 4 shots as traditional tabular baseline models do with 32 shots, as illustrated in Fig.~\ref{fig:compares}. Although STUNT and TabPFN show notable performance gains on individual datasets, their aggregate improvement over standard machine learning algorithms is modest. In addition, large language model-based methods—such as in-context learning, TABLET, and TabLLM—are unable to perform inference on tabular datasets characterized by a high number of features. By contrast, our framework incorporates both bagging and ensemble strategies and successfully handles datasets with extensive feature spaces, yielding strong outcomes. It should be emphasized that incorporating our bagging and ensemble methodology into other LLM-centered solutions is conceivable; however, this would necessitate doubling inference operations for each instance, substantially raising computational overhead and rendering it infeasible in practice.

\subsection{Analysis \& Discussion}
\label{sec:analysis}
\paragraph{Ablation study.} 
To assess the impact of key components, we conduct a series of ablation experiments by systematically removing individual elements: (1) \textsf{-Tuning}, where the linear model's weight tuning is skipped, using a simple sum over binary features for logit calculation; (2) \textsf{-Ensemble}, where the ensemble mechanism is excluded; (3) \textsf{-Description}, omitting feature descriptions; and (4) \textsf{-Reasoning}, eliminating the Step-1 process from reasoning instruction during rule generation.

Table~\ref{tab:ablation} reports the average change in AUC attributable to each ablated component, aggregated across datasets. Consistently, each removal leads to diminished performance. Notably, the positive effect of weight tuning intensifies with an increased number of shots, implying that as more training data becomes available, the framework is better able to accurately estimate rule significance, and thus tuning becomes increasingly valuable. Conversely, the importance of feature descriptions and reasoning instructions is most pronounced when only a small number of shots are provided, underscoring the vital role of leveraging the LLM's prior knowledge to enhance results in few-shot settings. Overall, the inclusion of our ensemble strategy reliably boosts model performance regardless of experimental configuration.

\vspace*{-0.12in}
\paragraph{Training \& inference time comparison.} We evaluate and contrast both training and inference durations across various approaches. These assessments utilize the Adult dataset with a training subset comprising 16 examples. For consistency, one A100 GPU serves as the standard computational resource, except for TabLLM, which relies on four A100 GPUs to enable model parallelism. As shown in Table~\ref{tab:cost}, the comparison of runtimes reveals that \textsf{FeatLLM} maintains a notably efficient inference speed, closely matching that of traditional baselines. In stark contrast, LLM-driven techniques such as In-context learning, TABLET, and TabLLM demonstrate significantly prolonged inference periods. Regarding training, \textsf{FeatLLM}'s time is similar to other LLM-based alternatives—it issues just 30 API requests, posing minimal cost concerns, and the process supports parallel execution.

\vspace*{-0.12in}
\paragraph{Quality of features generated by \textsf{FeatLLM}.}
To evaluate how informative the rule-based features generated by \textsf{FeatLLM} are for practical applications, we carry out a straightforward experiment. For each dataset, we compute the mean AUROC score between the derived features and the target labels and determine the proportion of features with AUROC scores exceeding 0.5. Results are compiled in Table~\ref{tab:quality}. The findings verify that most constructed rules pertain strongly to the intended prediction target and add significant value for the task at hand (see the “Before tuning” column). Furthermore, while fitting the linear model to the data, those rules contributing minimally (specifically, those associated with weights falling below the defined threshold) are automatically excluded (see the “After tuning” column). Here, we apply a threshold of 0.1, informed by the distribution of the model’s learned weights.

\vspace*{-0.12in}
\paragraph{Effect of adding spurious correlations.} To evaluate how different approaches manage extraneous spurious correlations, especially under limited labeled data, we designed an analysis. In this experiment, we employed the Heart dataset and artificially appended randomly selected columns from the Adult dataset, which bear no relevance to the Heart classification task. We assessed how model performance was affected as the number of such unrelated columns increased incrementally. To ensure an equitable evaluation, all methods utilized a fixed number of 8 labeled examples, omitting scenarios with an additional unlabeled dataset. As a result, approaches such as SCARF and STUNT were omitted from this assessment.

The effect of spurious correlations on model accuracy is depicted in Figure~\ref{fig:spurious}. When contrasted with baseline methods, \textsf{FeatLLM} demonstrated the most resilience, exhibiting minimal decline in performance due to spurious attributes. We attribute this robustness to the framework’s reliance on domain knowledge during rule extraction, which encourages learning the functional association between task and features and discourages rule formation from extraneous columns. As corroborated by our analysis, only a marginal fraction (1-2\%) of derived rules referenced the noisy attributes. However, it is important to note that when features from semantically similar datasets are integrated, \textsf{FeatLLM} is still susceptible to absorbing spurious relationships and can misconstrue the underlying causal links between labels and features (refer to Appendix~\ref{sec:spurious-appendix}).

\vspace*{-0.12in}
\paragraph{Hyper-parameter analysis}
\textsf{FeatLLM} employs two primary hyper-parameters: the total number of ensemble models used, and the quantity of rules extracted per class during each LLM invocation. To understand how these hyper-parameters influence model performance, we carried out systematic experiments. Our findings indicate that increasing the number of ensembles leads to improved outcomes, yet these gains plateau after approximately 20 ensembles (refer to Fig.~\ref{fig:appendix-hyperparameter1} in Appendix). With respect to the rule count, while we observed no statistically significant change in performance across various values, we selected 10 as the optimal number, as it provides a favorable balance between keeping the prompt concise and achieving strong results (see Fig.~\ref{fig:hyperparameter2}). We hypothesize that while increasing rule counts expands the input's dimensionality—and thus the potential information provided—it also intensifies the risk of overfitting, especially in few-shot learning contexts.

\section{Conclusion}
This paper introduces \textsf{FeatLLM}, a method that sets a new benchmark for few-shot tabular learning with LLMs. Rather than simply relying on LLMs to produce predictions for each instance, \textsf{FeatLLM} leverages LLMs for rule-based feature engineering that serves as decision criteria. Coupling rule induction with an ensemble strategy using bagging, \textsf{FeatLLM} achieves state-of-the-art performance while imposing minimal inference requirements and requiring only API-level interaction—eliminating the need for further model training or constraints due to data dimensionality. The model thus attains high accuracy for prediction tasks and presents a more effective approach for leveraging LLMs on tabular datasets, especially under data-efficient regimes.

Currently, \textsf{FeatLLM} targets only few-shot learning settings. In future work, we intend to broaden its scope to automate feature construction for datasets with more abundant samples, investigate feature forms beyond rules, and enhance interpretability—thereby facilitating its application to a broader set of machine learning challenges.

\section{Broader Impact}
The framework outlined in \textsf{FeatLLM} seeks to incorporate the extensive background knowledge of LLMs into tabular learning contexts, exceeding the performance ceilings of traditional supervised tabular approaches. By integrating LLM-derived prior knowledge, our research advances data-efficient learning and addresses prevalent issues like overfitting in scenarios where labeled examples are scarce. \textsf{FeatLLM} is especially useful for professionals in fields such as finance and healthcare, where data annotation is not only costly but also logistically difficult, and where LLMs’ pretrained domain expertise may prove invaluable. Looking ahead, a critical avenue for ensuring safe and responsible deployment involves thorough evaluation of societal biases and possible misinformation within the LLMs’ learned priors, as these may shape—or even supersede—the influence of available labeled data.

\end{document}